% ------------------------------------------------------------
% Page 90
% ------------------------------------------------------------

Aussi n’est-il pas étonnant que les camisards ne portent pas dans leurs cœurs M. de Ducros,
M. de Thomassy et quelques autres. Ils vont se venger, et entre les 11 et 13 mars 1703,
Castanet, « \textit{le camisard de l’Aigoual} » accompagné de La Rose, son fidèle lieutenant, et de
quelques autres au nombre d’une soixantaine pillèrent Connilergues. L’inventaire des effets
volés a été soigneusement noté par M. de Ducros qui, soulignant son état de Nouveau
Converti modèle, adresse une supplique à Monsieur l’Intendant pour obtenir un
remboursement, au moins partiel, de tout ce qui lui a été dérobé..

« \textit{vingt chemises, six cravates, quatre paires de manchettes, dix paires de draps fins pour le lit,
une veste d’écarlate avec les culottes de la même étoffe, quatre paires de bas, deux paires de
souliers une selle et une bride de cheval. Ils trouvèrent dans les culottes d’écarlate
mentionnées ci-dessus quatre-vingt-trois livres d’argent. Ils pénétrèrent dans la bibliothèque
et déchirèrent tous les livres de droit que le fils aîné du sieur de Campron, avocat, avait dans
son cabinet. Ils enlevèrent le miel de 12 ruches, du lard tout entier et la saucisse de trois
cochons. Ils eurent ou emportèrent deux charges de vin valant plus de 100 livres} ».*

Il souligne également qu’ils ont été obligés, sa famille et lui d’abandonner leur maison pour
se réfugier à Montpellier.

\vspace{0.2cm}

« \textit{Le 15 avril, les camisards faisant partie probablement de la même troupe de Castanet, écrit
Henri Bosc, revinrent chez le sieur Pierre du Cros, seigneur de Campron à Connilergues,
près de Meyrueis, où ils s’étaient déjà rendus au commencement du mois précédent. Ils
brûlèrent tout ce qu’ils trouvèrent et mirent le feu, en particulier, à cent soixante quintaux de
foin et à deux granges au lieu-dit « Les Oubretz ».} »

\vspace{0.2cm}

« \textit{Castanet, au cours du mois de mai 1703 devait revenir, une troisième fois piller le mas de
Connilergues. Ils étaient une quinzaine de camisards à cheval et un plus grand nombre à
pied. Ils brisèrent les vitres des appartements et s’emparèrent d’un veau, de deux des
meilleurs moutons du troupeau, d’une quarantaine de fromages et du miel de huit ruches
(chacun de ces fromages valait trois ou quatre livres) et du foin, foulé ou mangé par les 12 à
15 chevaux ou mulets qu’ils montaient} ».*

\vspace{0.2cm}

C’est en Novembre 1703 que les troupes royales, sous les ordres du maréchal de camp Julien,
se mirent à brûler les hameaux et villages de toute la région. Camprieu, Jontanels et Malbosc,
Campis et les Oubrets mais il n’est pas fait mention de Pourcarès ni de Connilergues,
épargnés peut-être, car appartenant à des familles de petite noblesse qui avaient abandonné la
Religion Prétendue Réformée pour se soumettre aux directives royales ?
Enfin vers 1750 la ferme de Connilergues, comme le Moulin d’Ayres d’ailleurs, apparaît sur
la carte de Cassini. Mais Connilergues est surmonté d’une bannière flottant au vent ; ce qui
indique l’importance de ce domaine appartenant à une famille de la noblesse.

\vspace{0.2cm}

A la date où je rédige ces quelques lignes, (Octobre 2007), on n’en sait pas plus et on ignore
l’histoire de Connilergues pendant plus de 100 ans.

\vspace{0.2cm}

« \textit{Les prédicants protestants des Cévennes et du bas Languedoc} » Ch. Bost 1912

« \textit{La guerre des Cévennes} » par Henri Bosc. I/ p.547, 567 et 713.

\begin{itemize}
    \item A.D. Hérault, C. 255
\end{itemize}
